\chapter{Pelvic Organ Prolapse Repair}

\section*{What Is This Surgery?}
Pelvic organ prolapse repair encompasses a diverse group of surgical procedures designed to correct the descent, or prolapse, of one or more pelvic organs into or through the vagina. This condition arises from weakened or damaged supportive tissues and ligaments of the pelvic floor, allowing organs such as the bladder (cystocele), rectum (rectocele), uterus, or the vaginal vault (after hysterectomy) to sag. The primary goal of these surgeries is to reposition the prolapsed organs, restore the structural integrity of the pelvic floor, and alleviate the associated bothersome symptoms, which can include a visible bulge, urinary dysfunction, bowel difficulties, and sexual discomfort. Surgical approaches vary widely, utilizing vaginal, laparoscopic, or open abdominal techniques, often with the aim of achieving durable anatomical correction and significant improvement in a patient's quality of life.

\section*{Alternative Names}
\begin{itemize}
  \item Anterior Colporrhaphy (for cystocele repair)
  \item Posterior Colporrhaphy (for rectocele repair)
  \item Sacrocolpopexy (abdominal or laparoscopic vaginal vault suspension)
  \item Sacrospinous Ligament Fixation (transvaginal apical suspension)
  \item Uterine Suspension or Hysteropexy (for uterine prolapse)
  \item Paravaginal Repair (for lateral compartment defects)
  \item Enterocele Repair (for small bowel prolapse)
\end{itemize}

\section*{Relevant Surgical Anatomy}
A thorough understanding of pelvic anatomy is paramount for successful pelvic organ prolapse repair. The pelvic floor is a complex structure comprising muscles, fascia, and ligaments that support the pelvic organs. Key muscular components include the levator ani muscle group (puborectalis, pubococcy and iliococcygeus muscles), which forms a sling-like diaphragm, and the perineal body, a fibromuscular structure central to the perineum. The endopelvic fascia, a connective tissue layer, provides crucial support, differentiating into the pubocervical fascia anteriorly (supporting the bladder), the rectovaginal fascia posteriorly (supporting the rectum), and the cardinal and uterosacral ligaments apically (supporting the uterus and vaginal vault).

Arterial supply to the pelvic floor and organs primarily originates from branches of the internal iliac artery, including the vaginal, uterine, middle rectal, and pudendal arteries. Venous drainage largely parallels the arterial supply, converging into the internal iliac veins. Nerve relations are critical; the pudendal nerve (S2-S4) innervates the external anal sphincter and perineal muscles, while autonomic nerves from the sacral plexus supply the pelvic organs. Lymphatic drainage follows the major blood vessels to the internal and external iliac lymph nodes. Surgical landmarks include the ischial spines (for sacrospinous ligament fixation), the sacral promontory (for sacrocolpopexy), and the arcus tendineus fascia pelvis (for paravaginal repair), all of which serve as critical attachment points for reconstructive efforts.

\section{Surgical Principle \& Rationale}
The fundamental principle of pelvic organ prolapse repair is to restore the anatomical support of the pelvic organs by addressing the underlying defects in the pelvic floor musculature and connective tissue. When these supports are stretched, torn, or weakened—often due to factors like childbirth trauma, chronic intra-abdominal pressure, or age-related tissue degradation—the pelvic organs descend. The operative concept involves identifying the specific compartment(s) of prolapse (anterior, posterior, or apical) and then either plicating (tightening) and reattaching native tissues, or augmenting these repairs with synthetic or biological grafts.

For anterior compartment prolapse (cystocele), the goal is to reinforce the pubocervical fascia and elevate the bladder neck. Posterior compartment prolapse (rectocele) requires plication of the rectovaginal fascia to support the rectum. Apical prolapse, involving the uterus or vaginal vault, necessitates suspension of the vaginal apex to strong, durable pelvic structures such as the sacral promontory (sacrocolpopexy) or the sacrospinous ligament (sacrospinous ligament fixation). The technical goals are multifaceted: achieving durable anatomical correction, preserving or restoring normal urinary, bowel, and sexual function, and minimizing the risk of recurrence or new onset symptoms. The choice of technique is guided by the specific type and severity of prolapse, patient comorbidities, and individual preferences, always aiming for the least invasive yet most effective solution.

\section{History \& Surgical Evolution}
Early attempts at pelvic organ prolapse repair date back to ancient times, but modern surgical techniques began to emerge in the 19th century. Anterior and posterior colporrhaphy, involving the plication of vaginal wall tissues, were established as foundational procedures by surgeons like Emmet and Kelly, and they remain cornerstones of vaginal prolapse repair today. The early 20th century saw the development of techniques to address apical prolapse, with the first descriptions of abdominal sacrocolpopexy by Huguier in 1859 and later popularized by Lane in 1962, offering a robust method for vaginal vault suspension.

The latter half of the 20th century brought innovations in transvaginal apical suspension, such as sacrospinous ligament fixation, which gained traction as a less invasive alternative to abdominal approaches. The early 2000s witnessed a significant shift with the widespread introduction of transvaginal mesh kits, marketed for their ease of use and perceived superior durability. However, a surge in reports of severe mesh-related complications, including erosion, chronic pain, and dyspareunia, led to regulatory actions by the FDA in the United States and a substantial decline in their use, particularly for anterior compartment repair. This era prompted a renewed focus on native-tissue repairs and a more cautious, evidence-based application of mesh. Concurrently, advancements in minimally invasive surgery have led to the increasing adoption of laparoscopic and robotic sacrocolpopexy, offering the benefits of reduced recovery time with excellent long-term outcomes.

\section{Clinical Indications}
Pelvic organ prolapse repair is indicated for women experiencing symptomatic pelvic organ prolapse that significantly impacts their quality of life and has not responded adequately to conservative management.

\begin{itemize}
  \item A visible or palpable bulge at the vaginal opening is the most common and compelling reason for surgical intervention.
  \item A persistent sensation of pelvic pressure, heaviness, or a feeling that organs are "falling out" despite conservative measures.
  \item Significant urinary symptoms, including a weak stream, incomplete bladder emptying, recurrent urinary tract infections, or stress urinary incontinence exacerbated by prolapse.
  \item Troublesome bowel symptoms such as chronic constipation, fecal urgency, difficulty with defecation, or the need for manual splinting of the perineum to facilitate bowel movements.
  \item Sexual dysfunction, including dyspareunia (painful intercourse) or a sensation of vaginal looseness, directly attributable to the prolapse.
  \item Prolapse that interferes with daily activities, exercise, or work, leading to a substantial reduction in a woman's overall functional capacity.
  \item Failure of, intolerance to, or patient preference against non-surgical treatments such as pelvic floor physiotherapy or vaginal pessaries.
  \item Progressive anatomical prolapse, even if symptoms are initially mild, particularly in younger women where long-term support is a concern.
\end{itemize}

\section{Clinical Positioning}
Pelvic organ prolapse repair is generally considered a secondary treatment option in the overall management pathway for symptomatic prolapse. Initial management typically involves conservative strategies, which are often highly effective for many women. These include lifestyle modifications (e.g., weight loss, management of chronic cough or constipation), pelvic floor muscle training (pelvic floor physiotherapy), and the use of vaginal pessaries. A trial of a pessary is often recommended as a first-line intervention, as it can provide significant symptom relief and help patients determine if anatomical correction will resolve their specific complaints.

Surgery is offered when conservative measures fail to adequately alleviate symptoms, are not tolerated by the patient, or are declined due to personal preference or lifestyle considerations. In this context, surgery serves as the primary definitive treatment once the decision to operate has been made. The specific surgical technique chosen (e.g., vaginal vs. abdominal, native tissue vs. mesh augmentation) is highly individualized, based on the type and severity of prolapse, patient age, comorbidities, sexual activity, desire for future childbearing, and surgeon expertise. While surgery is not typically the first step, it is the most effective and durable solution for many women with bothersome, advanced prolapse.

\section{Surgical Technique \& Procedure}
Pelvic organ prolapse repair is performed under general or regional (spinal/epidural) anesthesia, with patient positioning tailored to the surgical approach. For vaginal repairs, the patient is placed in the dorsal lithotomy position. For abdominal or laparoscopic approaches, the patient is typically supine with Trendelenburg tilt.

The specific steps vary significantly by the type of prolapse and chosen technique:

\begin{enumerate}
  \item  **Anesthesia and Positioning**: General or regional anesthesia is administered. The patient is positioned in dorsal lithotomy for vaginal approaches, or supine with appropriate padding for abdominal/laparoscopic cases. The surgical field is prepped and draped.
  \item  **Incision**: For vaginal repairs, an incision is made in the anterior or posterior vaginal wall, or circumferentially around the cervix/vaginal apex. For abdominal sacrocolpopexy, a transverse lower abdominal incision (open) or multiple small port incisions (laparoscopic/robotic) are made.
  3.  **Dissection and Identification of Defects**: The vaginal wall is carefully dissected from the underlying bladder (anterior repair) or rectum (posterior repair). For apical prolapse, the vaginal apex is identified, and surrounding structures like the uterosacral ligaments or sacrospinous ligaments are exposed. In abdominal sacrocolpopexy, the peritoneum is opened to expose the sacral promontory and the vaginal apex.
  4.  **Repair of Anterior Compartment (Anterior Colporrhaphy)**: The pubocervical fascia is identified, plicated with absorbable sutures to reinforce the bladder support, and the bladder is elevated. Excess vaginal skin may be trimmed.
  5.  **Repair of Posterior Compartment (Posterior Colporrhaphy)**: The rectovaginal fascia is identified and plicated with absorbable sutures to support the rectum and reconstruct the perineal body. Excess vaginal skin may be trimmed.
  6.  **Apical Suspension (e.g., Sacrocolpopexy or Sacrospinous Ligament Fixation)**:
      \begin{itemize}
        \item **Sacrocolpopexy**: A synthetic mesh graft (typically polypropylene) is attached to the vaginal apex (or posterior cervix/uterus for hysteropexy) and then secured to the anterior longitudinal ligament over the sacral promontory. This can be performed via open, laparoscopic, or robotic approaches.
        \item **Sacrospinous Ligament Fixation**: Through a vaginal incision, the sacrospinous ligament is identified, and sutures are placed through it and then secured to the vaginal apex to suspend it.
      \end{itemize}
  7.  **Concomitant Procedures**: If stress urinary incontinence is present, a mid-urethral sling procedure may be performed concurrently. If a hysterectomy is indicated (e.g., for uterine prolapse or other gynecological pathology), it is performed before apical suspension.
  8.  **Closure**: The vaginal incisions are closed with absorbable sutures. Abdominal incisions are closed in layers. A vaginal pack may be inserted temporarily, and a Foley catheter is often placed to drain the bladder.
\end{enumerate}
The duration of surgery typically ranges from 1 to 3 hours, depending on the complexity and number of compartments repaired. Patients undergoing vaginal repairs may be discharged the same day or after a short overnight stay, while more extensive abdominal procedures may require 1-2 nights in the hospital.

\section*{Preoperative Workup \& Preparation}
A comprehensive preoperative workup is essential to optimize outcomes and minimize risks for pelvic organ prolapse repair. This begins with a detailed medical history, including parity, obstetric history, prior surgeries, and a thorough review of urinary, bowel, and sexual symptoms. A physical examination, including a standardized pelvic examination (e.g., using the POP-Q system - Pelvic Organ Prolapse Quantification), is performed to accurately assess the type and severity of prolapse in all compartments.

\begin{itemize}
  \item **Medical Evaluation**: A general medical assessment, including cardiac and pulmonary evaluation, is conducted to ensure the patient is fit for anesthesia and surgery. This may involve an ECG, chest X-ray, and blood tests (complete blood count, electrolytes, renal function, coagulation profile).
  \item **Urodynamic Studies**: If significant urinary symptoms, such as stress urinary incontinence or voiding dysfunction, are present, multichannel urodynamic testing may be performed to characterize bladder function and guide surgical planning.
  \item **Imaging**: While not routinely required, imaging such as pelvic ultrasound or MRI may be used in complex cases to delineate anatomy, rule out other pelvic pathology, or assess the extent of prolapse.
  \item **Bowel Preparation**: For posterior repairs or procedures involving the rectovaginal space, a bowel preparation (e.g., clear liquid diet and oral laxatives) may be advised to reduce the risk of infection and improve surgical field visualization.
  \item **Medication Adjustment**: Patients are instructed to discontinue blood-thinning medications (e.g., aspirin, NSAIDs, warfarin, novel oral anticoagulants) for a specified period before surgery, as advised by the surgical team.
  \item **NPO Status**: Patients must be nil per os (NPO) for at least 6-8 hours prior to surgery to reduce the risk of aspiration during anesthesia.
  \item **Prophylactic Antibiotics**: Intravenous broad-spectrum antibiotics are typically administered within 60 minutes prior to incision to reduce the risk of surgical site infection.
  \item **Informed Consent**: A detailed discussion of the proposed procedure, alternative treatments, potential benefits, and specific risks (including recurrence, mesh complications if applicable, and impact on sexual function) is conducted, and informed consent is obtained.
  \item **Smoking Cessation**: Patients are strongly advised to stop smoking several weeks prior to surgery to improve wound healing and reduce respiratory complications.
\end{itemize}

\section*{Contraindications \& Precautions}
Careful consideration of contraindications and precautions is crucial to ensure patient safety and optimize surgical outcomes.

\textbf{Absolute Contraindications:}
\begin{itemize}
  \item Active pelvic infection (e.g., pelvic inflammatory disease, vaginitis, urinary tract infection) must be treated and resolved prior to surgery.
  \item Uncontrolled systemic medical conditions (e.g., severe heart failure, unstable angina, uncontrolled diabetes, severe coagulopathy) that pose an unacceptable anesthetic or surgical risk.
  \item Pregnancy.
  \item Active malignancy in the pelvic region that requires primary oncological treatment.
\end{itemize}

\textbf{Relative Contraindications \& Precautions:}
\begin{itemize}
  \item **Desire for Future Childbearing**: Women who plan to have more children are generally advised to delay surgery, as subsequent pregnancy and vaginal delivery can disrupt the repair and lead to recurrence. Uterine-sparing procedures (hysteropexy) may be considered in selected cases.
  \item **Use of Synthetic Mesh**: While mesh can provide durable support, its use requires careful consideration due to the risk of erosion, infection, chronic pain, and dyspareunia. Many centers restrict mesh use to specific high-risk situations, such as recurrent prolapse or severe apical defects, and avoid it for routine anterior or posterior repairs.
  \item **Significant Obesity**: While not an absolute contraindication, severe obesity increases surgical complexity, anesthetic risks, and may impair the long-term durability of the repair due to increased intra-abdominal pressure. Weight loss is often recommended preoperatively.
  \item **Chronic Conditions Increasing Intra-abdominal Pressure**: Conditions such as chronic cough (e.g., from COPD, asthma) or chronic constipation can place excessive strain on the pelvic floor, increasing the risk of recurrence. These conditions should be optimally managed before and after surgery.
  \item **Smoking**: Smoking impairs wound healing and increases the risk of postoperative complications. Smoking cessation is strongly encouraged.
  \item **Unrealistic Patient Expectations**: Patients must have a clear understanding of the potential outcomes, risks, and the possibility of recurrence. Thorough counseling is essential.
  \item **Prior Pelvic Radiation**: Radiation therapy can compromise tissue quality and vascularity, increasing the risk of complications and affecting the durability of repairs.
\end{itemize}

\section*{Complications \& Risks}
As with any surgical procedure, pelvic organ prolapse repair carries both general surgical risks and specific procedure-related complications.

\textbf{General Surgical Complications:}
\begin{itemize}
  \item **Bleeding**: Intraoperative or postoperative hemorrhage, potentially requiring blood transfusion or re-operation.
  \item **Infection**: Surgical site infection (vaginal, abdominal wound), urinary tract infection, or more severe pelvic infection.
  \item **Anesthesia Risks**: Adverse reactions to anesthetic agents, respiratory complications (e.g., pneumonia, atelectasis), cardiovascular events (e.g., myocardial infarction, stroke).
  \item **Thromboembolic Events**: Deep vein thrombosis (DVT) or pulmonary embolism (PE), despite prophylactic measures.
  \item **Pain**: Postoperative pain, which is usually manageable with analgesics but can occasionally be chronic.
\end{itemize}

\textbf{Procedure-Specific Complications:}
\begin{itemize}
  \item **Recurrence of Prolapse**: The most common long-term complication, with rates varying significantly based on the type of prolapse, surgical technique, and patient factors. Recurrence can involve the same or a different compartment.
  \item **Injury to Adjacent Organs**:
    \begin{itemize}
      \item **Bladder Injury**: During anterior dissection or suture placement, potentially leading to fistula formation (vesicovaginal fistula).
      \item **Ureteric Injury**: Particularly during apical suspension or hysterectomy, potentially leading to hydronephrosis or fistula (ureterovaginal fistula).
      \item **Bowel Injury**: During posterior dissection or abdominal approaches, potentially leading to peritonitis or fistula (rectovaginal fistula).
    \end{itemize}
  \item **New or Worsening Urinary Symptoms**:
    \begin{itemize}
      \item **De novo Stress Urinary Incontinence**: Unmasking of occult SUI after prolapse reduction.
      \item **Urinary Urgency or Frequency**: New onset or worsening of overactive bladder symptoms.
      \item **Voiding Dysfunction**: Difficulty emptying the bladder, sometimes requiring temporary or permanent catheterization.
    \end{itemize}
  \item **Mesh-Related Complications (if mesh is used)**:
    \begin{itemize}
      \item **Mesh Erosion/Exposure**: The mesh material can erode through the vaginal wall or into adjacent organs (bladder, rectum).
      \item **Chronic Pain**: Persistent pelvic, vaginal, or groin pain, sometimes severe and debilitating.
      \item **Dyspareunia**: Painful intercourse, which can be new onset or worsened.
      \item **Mesh Contraction/Shrinkage**: Leading to vaginal shortening or narrowing.
    \end{itemize}
  \item **Dyspareunia**: New onset or worsening of painful intercourse, even with native tissue repairs, due to vaginal shortening, narrowing, or scarring.
  \item **Nerve Injury**: Rare, but can occur, such as pudendal nerve injury leading to perineal numbness or pain.
  \item **Fistula Formation**: Abnormal communication between the vagina and bladder (vesicovaginal), rectum (rectovaginal), or ureter (ureterovaginal), a rare but serious complication requiring further intervention.
\end{itemize}

\section*{Postoperative Recovery Timeline}
The postoperative recovery timeline for pelvic organ prolapse repair varies depending on the extent of the surgery and the individual patient, but generally follows a predictable course.

Immediately after surgery, patients are monitored in the Post-Anesthesia Care Unit (PACU) for recovery from anesthesia, pain control, and vital sign stability. A Foley catheter may remain in place for a few hours or days, especially if bladder function was compromised or if a concomitant anti-incontinence procedure was performed, to ensure adequate bladder drainage and allow for initial healing. Oral analgesics are typically prescribed for pain management, and early mobilization is encouraged to reduce the risk of deep vein thrombosis and promote bowel function. Most patients are able to tolerate a light diet within hours of surgery.

During the first 1-2 weeks, patients are advised to avoid heavy lifting (typically >10 lbs), strenuous exercise, and prolonged standing to protect the surgical repair. Constipation is a common concern and is managed with increased fluid intake, a high-fiber diet, and stool softeners to prevent straining, which can put undue pressure on the healing pelvic floor. Vaginal discharge or light bleeding is common. Most women can resume light daily activities within 2-3 weeks. Full recovery, including the resumption of strenuous activities and sexual intercourse, is typically advised after 6-8 weeks, once tissue healing is confirmed by the surgeon. Patients are educated on warning signs such as fever, increasing pain, heavy bleeding, or difficulty urinating, which warrant immediate medical attention.

\section*{Surgical Findings \& Pathology}
During pelvic organ prolapse repair, the surgeon meticulously documents the specific anatomical defects identified. This includes the degree and type of prolapse in each compartment (e.g., cystocele, rectocele, enterocele, uterine/vault prolapse), the integrity of the endopelvic fascia, the condition of the levator ani muscles, and the presence of any associated conditions like stress urinary incontinence. The surgeon notes the specific tissues plicated, the ligaments utilized for suspension, and whether any graft material (synthetic mesh or biological graft) was used, including its type and fixation points. Any intraoperative complications, such as organ injury or excessive bleeding, are also thoroughly documented.

Pathology reports on resected tissues provide crucial histological confirmation. If a hysterectomy is performed, the uterus and cervix are sent for routine pathological examination to rule out any underlying uterine or cervical pathology. In cases of anterior or posterior colporrhaphy, excess vaginal wall tissue may be submitted, which typically shows benign squamous epithelium with underlying fibromuscular tissue, often with signs of atrophy or chronic inflammation. If mesh is removed due to complications, the pathologist examines the mesh material and surrounding tissue for signs of inflammation, foreign body reaction, infection, or necrosis. These findings help correlate clinical presentation with microscopic changes and guide future management.

\section{Expected Postoperative Course}
A normal and successful postoperative course following pelvic organ prolapse repair is characterized by a progressive resolution of the preoperative bulge and pressure symptoms, along with significant improvement in urinary, bowel, and sexual function. Patients should experience a gradual decrease in surgical pain, which is typically well-controlled with oral analgesics and resolves within a few weeks. The vaginal incision site heals without complication, and any temporary vaginal discharge or spotting subsides. Bladder function should normalize, with patients able to void spontaneously and completely, and any new or worsened urinary symptoms should be transient. Bowel movements should resume regularly without straining. Ultimately, the expectation is a return to normal daily activities, exercise, and sexual function without discomfort or recurrence of prolapse symptoms, leading to a substantial improvement in overall quality of life.

\section{Postoperative Care \& Follow-up}
Postoperative care and follow-up are critical for ensuring optimal healing, monitoring for complications, and assessing the long-term durability of the repair.

\begin{itemize}
  \item **Initial Follow-up**: The first postoperative visit typically occurs 2-4 weeks after surgery to assess wound healing, remove any non-absorbable sutures if present, and address any immediate concerns.
  \item **Activity Restrictions**: Patients are advised to avoid heavy lifting, strenuous exercise, and sexual intercourse for at least 6-8 weeks to allow for adequate tissue healing and integration of any graft material.
  \item **Bowel and Bladder Management**: Continued emphasis on preventing constipation through diet, fluids, and stool softeners is crucial. Patients are monitored for complete bladder emptying.
  \item **Pelvic Floor Physiotherapy**: Referral to a pelvic floor physiotherapist may be recommended to strengthen pelvic floor muscles, improve core stability, and optimize long-term support.
  \item **Long-term Follow-up**: Subsequent visits are usually scheduled at 6 months and 1 year, and then annually, to assess for recurrence of prolapse, new onset symptoms, and overall patient satisfaction.
  \item **Warning Signs**: Patients are educated to contact their surgeon immediately if they experience fever, increasing pelvic pain, heavy vaginal bleeding, foul-smelling discharge, difficulty urinating, or a new vaginal bulge.
\item **Lifestyle Modifications**: Maintenance of a healthy weight, avoidance of chronic straining (e.g., with constipation or cough), and smoking cessation are encouraged to support the long-term success of the repair.
\end{itemize}

\section*{Epidemiology}
Pelvic organ prolapse is a highly prevalent condition, with epidemiological studies estimating that up to 50\% of parous women exhibit some degree of prolapse on physical examination, although only a fraction of these women are symptomatic enough to seek treatment. The incidence of symptomatic prolapse increases significantly with age, particularly after menopause, due to hormonal changes affecting tissue integrity. The lifetime risk of undergoing surgery for pelvic organ prolapse or stress urinary incontinence is substantial, generally estimated to be between 10\% and 20\%.

Risk factors for pelvic organ prolapse are well-established and include vaginal childbirth (especially multiple, traumatic, or instrumented deliveries), increasing age, obesity (which increases intra-abdominal pressure), chronic conditions that lead to repetitive straining (e.g., chronic constipation, chronic cough), prior hysterectomy (which removes apical support), and a family history suggesting a genetic predisposition to connective tissue weakness. As global populations age, the number of women affected by pelvic organ prolapse and requiring surgical intervention is projected to rise, posing a significant public health challenge. While not directly associated with mortality, pelvic organ prolapse carries significant morbidity, severely impacting a woman's quality of life, physical activity, and psychological well-being.

\section*{Outcomes \& Prognosis}
The outcomes and prognosis following pelvic organ prolapse repair are generally favorable, with the vast majority of women experiencing significant improvement in their symptoms and quality of life. Success rates, defined by anatomical correction and symptom resolution, vary depending on the type of prolapse, the surgical technique employed, and patient-specific factors. For apical suspension procedures like abdominal or laparoscopic sacrocolpopexy, success rates are high, often exceeding 85-90\% at 1-2 years and remaining durable over longer follow-up periods, as demonstrated by landmark trials such as the CARE (Colpopexy and Advanced Pelvic Surgery) trial. Native-tissue anterior and posterior colporrhaphy, while effective, tend to have higher anatomical recurrence rates, particularly for anterior compartment prolapse, though symptom improvement can still be substantial.

Recurrence of prolapse, either in the same or a different compartment, remains the most common long-term problem, affecting 15-30\% of women over 5 years, and some women may require a second or even third procedure. Prognostic factors influencing outcomes include the severity of the initial prolapse, patient age, body mass index, presence of connective tissue disorders, and adherence to postoperative lifestyle modifications. Despite the risk of recurrence, overall patient satisfaction with prolapse surgery is high, with studies consistently reporting significant improvements in pelvic floor symptoms, sexual function, and overall quality of life. When surgery is appropriately indicated and meticulously performed, durable anatomical correction and functional improvement are the expected outcomes.

\section{References}
\begin{enumerate}
  \item MedlinePlus. Pelvic organ prolapse surgery. U.S. National Library of Medicine; 2024. Available from: https://medlineplus.gov/pelvicorganprolapsesurgery.html
  \item Hoffman BL, Schorge JO, Bradshaw KD, Halvorson LM, Schaffer JI, Corton MC. Williams Gynecology. 4th ed. New York, NY: McGraw-Hill Education; 2020.
  \item American College of Obstetricians and Gynecologists. ACOG Practice Bulletin No. 220: Pelvic Organ Prolapse. Obstet Gynecol. 2020;136(3):e66-e83.
  \item Olsen AL, Smith VJ, Bergstrom J, Colling JC, Clark RA. Epidemiology of surgically managed pelvic organ prolapse and urinary incontinence. Obstet Gynecol. 1997;89(4):501-506.
\end{enumerate}

\clearpage